\documentclass[11pt,a4paper]{article}
\usepackage{fontspec}
\usepackage{polyglossia}
\setdefaultlanguage{french}
\setmainfont{Latin Modern Roman}
\setmonofont{Latin Modern Mono}
\usepackage[margin=2.5cm]{geometry}
\usepackage{graphicx,booktabs,tabularx,array,amsmath,amssymb,xcolor,enumitem,float,caption,fancyhdr,hyperref}
\hypersetup{colorlinks=true,linkcolor=blue!50!black,urlcolor=blue!50!black,citecolor=blue!50!black}
\captionsetup{labelfont=bf,font=small}
\setlength{\parindent}{0pt}\setlength{\parskip}{5pt}
\renewcommand{\arraystretch}{1.2}
\pagestyle{fancy}\fancyhf{}
\fancyhead[L]{\small ENCG Settat -- BD \& DS}
\fancyhead[R]{\small TAISA WARDA -- 25007836}
\fancyfoot[C]{\thepage}

\begin{document}
%==================== PAGE DE GARDE ====================
\thispagestyle{empty}
\begin{center}
\includegraphics[height=2.2cm]{logo_universite_hassan1.png}\hfill
\includegraphics[height=2.2cm]{logo_encg_settat.png}\\[1.2cm]
{\large Université Hassan 1\textsuperscript{er} de Settat}\\[2pt]
{\Large\bfseries École Nationale de Commerce et de Gestion de Settat}\\[2pt]
ENCG Settat\\[1cm]
\rule{\textwidth}{0.6pt}\\[6pt]
Cours BD \& DS : examen individuel (Partie B)\\[2pt]
{\large Rapport de projet de science des données appliquée au commerce international}\\[10pt]
{\LARGE\bfseries Prévision des exportations marocaines}\\[4pt]
{\large Régression par forêt aléatoire sur les données UN Comtrade}\\[6pt]
\rule{\textwidth}{0.6pt}
\end{center}
\vspace{1.2cm}
\noindent
\begin{minipage}[c]{0.64\textwidth}
\renewcommand{\arraystretch}{1.5}
\begin{tabular}{@{}lp{6.2cm}@{}}
\textbf{Réalisé par} & \textbf{TAISA WARDA}\\
\textbf{Code Apogée} & \textbf{25007836}\\
\textbf{Groupe et filière} & \textbf{Groupe S7 CI}\\
\textbf{Sujet} & \textbf{CI01 -- Prévision des exportations}\\
\textbf{Algorithme principal} & \textbf{n°12 -- Random Forest Regressor}\\
\textbf{Enseignant} & \textbf{A. Larhlimi}\\
\end{tabular}
\end{minipage}\hfill
\begin{minipage}[c]{0.30\textwidth}
\centering
\includegraphics[width=\textwidth]{photo_etudiant.jpg}
\end{minipage}

\vspace{1cm}
\textbf{Dépôt GitHub :} \url{https://github.com/wardataisa/25007836-Taisa-CI.git}

\vfill
\begin{center}
\textbf{Année universitaire 2026-2027}\\
\textbf{Remis le 26 novembre 2026}
\end{center}
\newpage
\tableofcontents
\newpage

%==================== 1 INTRODUCTION ====================
\section{Introduction}
\textbf{Contexte.} Le commerce extérieur est un moteur de la croissance du Maroc, soutenu par des accords de libre-échange et des stratégies sectorielles (automobile, agriculture, phosphates). Anticiper la valeur des exportations aide les décideurs publics et les entreprises à prévoir les rentrées de devises, à ajuster la politique commerciale et à consolider la balance commerciale dans un contexte mondial volatil.

\textbf{Problématique.} Peut-on prédire la valeur des exportations marocaines (\texttt{primaryValue}) à partir du pays partenaire, de la catégorie de produit (\texttt{cmdCode}), de l'année et du flux commercial ?

\textbf{Objectif mesurable.} Fixé avant l'expérimentation : obtenir un coefficient de détermination $R^2>0{,}65$ sur le jeu de test et réduire la MAE d'au moins 30\,\% par rapport à la référence naïve (moyenne du jeu d'apprentissage).

\textbf{Nature du problème.} Régression supervisée : la cible est une valeur monétaire continue en dollars américains.

\textbf{Organisation du rapport.} La section~\ref{sec:eda} présente l'état de l'art et la section~\ref{sec:liés} les travaux liés. Les sections~\ref{sec:data} à~\ref{sec:proto} décrivent les données, leur exploration, leur préparation, l'algorithme retenu et le protocole expérimental. La section~\ref{sec:res} présente les résultats, discutés en section~\ref{sec:disc}, avant la conclusion (section~\ref{sec:concl}).

%==================== 2 ETAT DE L'ART ====================
\section{État de l'art}\label{sec:eda}
La prévision des flux commerciaux repose historiquement sur le modèle de gravité, qui explique les échanges bilatéraux par la taille économique des partenaires et la distance qui les sépare \cite{tinbergen,anderson}. Les approches d'apprentissage automatique, notamment les méthodes d'ensemble d'arbres, permettent de capter des non-linéarités et des interactions sans spécification préalable \cite{breiman}.

\begin{table}[H]\centering\small
\caption{Panorama des approches pour la prévision des flux commerciaux}
\begin{tabularx}{\textwidth}{@{}p{2.6cm}XXXc@{}}
\toprule
Approche & Principe & Atouts & Limites & Réf.\\\midrule
Référence naïve & Prévoir la moyenne (ou la dernière valeur) & Simple, sans paramètre & Ignore partenaires et produits & ---\\
Modèle de gravité (approche classique) & Régression log-linéaire sur PIB, distance, frontières & Interprétable, fondé en théorie & Exige des variables macro ; linéaire & \cite{tinbergen,anderson}\\
Apprentissage automatique (forêt aléatoire) & Moyenne de nombreux arbres de décision décorrélés & Non-linéarités, interactions, robuste au bruit & Extrapole mal hors de la plage observée & \cite{breiman}\\
\bottomrule
\end{tabularx}
\end{table}

%==================== 3 TRAVAUX LIES ====================
\section{Travaux liés}\label{sec:liés}
\begin{table}[H]\centering\small
\caption{Travaux liés}
\begin{tabularx}{\textwidth}{@{}p{2.6cm}p{2.3cm}p{2.3cm}XX@{}}
\toprule
Travail & Données & Méthode & Résultat principal & Différence avec ce travail\\\midrule
Anderson \& van Wincoop, 2003 \cite{anderson} & Flux bilatéraux États-Unis--Canada et entre pays & Gravité structurelle & Les frontières réduisent fortement les échanges ; la « résistance multilatérale » doit être contrôlée & Ici : pas de variables macro, modèle non paramétrique\\
Yotov et al., 2016 \cite{yotov} & Flux bilatéraux mondiaux & Gravité structurelle (PPML) & Guide d'estimation de référence pour l'analyse des politiques commerciales & Ici : objectif de prévision, pas d'estimation causale\\
\bottomrule
\end{tabularx}
\end{table}

\textbf{Positionnement.} Ce travail reprend l'idée que partenaire, produit et période structurent les échanges, mais l'aborde sous l'angle prédictif avec une forêt aléatoire. Il n'intègre pas les déterminants macroéconomiques des modèles de gravité, ce qui constitue sa principale limite.

%==================== 4 DONNEES ====================
\section{Jeu de données}\label{sec:data}
\begin{itemize}[leftmargin=1.2em]
\item \textbf{Nom et source :} base UN Comtrade (statistiques officielles du commerce de marchandises), Nations unies \cite{comtrade}, \url{https://comtradeplus.un.org/}.
\item \textbf{Date d'extraction :} 19/11/2026.
\item \textbf{Licence :} conditions d'utilisation publiques de UN Comtrade.
\item \textbf{Taille :} 8\,288 observations et 47 variables, fréquence annuelle, de 2016 à 2025 d'après la figure~\ref{fig:evol} (l'en-tête du notebook indique 2016--2024).
\item \textbf{Unité d'observation :} une ligne représente le flux annuel de marchandises du Maroc (déclarant, code 504) vers un partenaire pour un produit du Système Harmonisé.
\item \textbf{Devise :} dollars américains (USD).
\item \textbf{Variable cible :} \texttt{primaryValue}, valeur de la transaction.
\item \textbf{Adéquation :} le jeu contient les quatre dimensions de la problématique (année, flux, partenaire, produit) et la valeur à prédire.
\item \textbf{Limites :} aucune variable macroéconomique ; présence de la ligne partenaire « World » (agrégat) à côté des pays ; période courte (dix ans) marquée par des chocs (pandémie, inflation).
\end{itemize}

\begin{table}[H]\centering\small
\caption{Dictionnaire des principales variables}
\begin{tabularx}{\textwidth}{@{}lXlllc@{}}
\toprule
Variable & Description & Type & Unité & Rôle & Manq.\,(\%)\\\midrule
\texttt{refYear} & Année de la transaction & entier & année & explicative & 0\\
\texttt{flowDesc} & Sens du flux (Export) & catégoriel & --- & explicative & 0\\
\texttt{partnerDesc} & Pays partenaire & catégoriel & --- & explicative & 0\\
\texttt{cmdCode} & Code produit (Système Harmonisé) & entier & code SH & explicative & 0\\
\texttt{primaryValue} & Valeur de la transaction & numérique & USD & cible & 0\\
\bottomrule
\end{tabularx}
\end{table}

%==================== 5 EXPLORATION ====================
\section{Exploration}\label{sec:expl}
\textbf{Statistiques descriptives.} Le jeu compte 8\,288 lignes. La valeur des transactions couvre plus de neuf ordres de grandeur, de quelques dollars à plusieurs dizaines de milliards (figure~\ref{fig:dist}).

\textbf{Qualité des données.} Le notebook retire les lignes dont \texttt{primaryValue}, \texttt{refYear} ou \texttt{flowDesc} est manquant ; les colonnes retenues n'ont pas de valeurs manquantes. Les valeurs extrêmes ont été conservées car, en commerce international, elles correspondent à de vraies transactions de grande taille.

\begin{figure}[H]\centering
\includegraphics[width=0.85\textwidth]{figures/fig1_evolution.png}
\caption{Évolution des exportations marocaines par année (milliards USD)}\label{fig:evol}
\end{figure}
\textbf{Lecture de la figure~\ref{fig:evol}.} Les montants cumulés progressent de façon régulière entre 2016 et 2023, puis grimpent nettement en 2024 avant de baisser en 2025. Ce pic est à interpréter avec prudence : les sommes mélangent la ligne « World » et les lignes par pays, ce qui peut compter plusieurs fois les mêmes flux.

\begin{figure}[H]\centering
\includegraphics[width=0.85\textwidth]{figures/fig2_partenaires.png}
\caption{Dix premiers partenaires commerciaux (milliards USD)}\label{fig:part}
\end{figure}
\textbf{Lecture de la figure~\ref{fig:part}.} « World » domine car c'est un agrégat. Parmi les pays, l'Espagne et la France ressortent nettement, devant l'Italie, l'Allemagne et le Brésil : les exportations sont concentrées sur quelques partenaires, surtout européens.

\begin{figure}[H]\centering
\includegraphics[width=0.85\textwidth]{figures/fig3_distribution.png}
\caption{Distribution de \texttt{primaryValue} (échelle logarithmique)}\label{fig:dist}
\end{figure}
\textbf{Lecture de la figure~\ref{fig:dist}.} La distribution est approximativement symétrique en échelle log, centrée autour de $10^5$--$10^7$ USD, avec une longue queue de très grosses transactions. Cette asymétrie explique que l'erreur en dollars soit dominée par quelques observations.

%==================== 6 PREPARATION ====================
\section{Préparation}\label{sec:prep}
\begin{itemize}[leftmargin=1.2em]
\item \textbf{Nettoyage :} suppression des lignes sans valeur pour la cible, l'année ou le flux.
\item \textbf{Variables retenues :} \texttt{refYear}, \texttt{flowDesc}, \texttt{partnerDesc}, \texttt{cmdCode}. Aucune variable dérivée n'a été créée.
\item \textbf{Encodage et mise à l'échelle :} \texttt{StandardScaler} sur \texttt{refYear} et \texttt{cmdCode} ; \texttt{OneHotEncoder} (\texttt{handle\_unknown='ignore'}) sur \texttt{flowDesc} et \texttt{partnerDesc}. La standardisation n'est pas nécessaire pour une forêt aléatoire, mais elle est inoffensive.
\end{itemize}

\textbf{Prévention des fuites.}
\begin{itemize}[label=$\square$,leftmargin=1.5em]
\item Transformations ajustées sur l'apprentissage seulement, dans un \texttt{Pipeline}.
\item Découpage chronologique, sans mélange : lignes triées par \texttt{refYear}, 80\,\% premières pour l'apprentissage.
\item Aucune variable ne contient la cible : \texttt{cifvalue} et \texttt{fobvalue}, proches de \texttt{primaryValue}, ne sont pas utilisées.
\item Jeu de test évalué une seule fois, à la fin.
\end{itemize}
Le découpage par position sur les lignes triées par année peut couper une année en deux, car plusieurs lignes partagent la même année.

%==================== 7 ALGORITHME ====================
\section{Algorithme}\label{sec:algo}
\textbf{Choix.} Algorithme n°12 : forêt aléatoire de régression (\texttt{RandomForestRegressor}, scikit-learn \cite{sklearn}). Elle capte les relations non linéaires et les interactions entre produits et partenaires, en accord avec l'état de l'art \cite{breiman}. Une régression linéaire a été écartée car elle ne représente pas l'effet combiné produit $\times$ pays.

\textbf{Principe.} On entraîne $B$ arbres de décision sur des échantillons bootstrap des données, chacun avec un sous-ensemble aléatoire de variables à chaque coupure. La prédiction est la moyenne des arbres.

\textbf{Formule essentielle.}
\begin{equation}
\hat y=\frac{1}{B}\sum_{b=1}^{B}T_b(x)
\end{equation}
où $\hat y$ est la valeur prédite, $B$ le nombre d'arbres, $x$ le vecteur de caractéristiques de l'observation et $T_b(x)$ la prédiction de l'arbre $b$.

\textbf{Paramètres appris.} Structure des arbres, seuils de coupure et valeurs moyennes aux feuilles.

\begin{table}[H]\centering\small
\caption{Hyperparamètres}
\begin{tabularx}{\textwidth}{@{}lXll@{}}
\toprule
Hyperparamètre & Rôle et effet attendu & Valeurs envisagées & Valeur retenue\\\midrule
\texttt{n\_estimators} & Nombre d'arbres : plus il est grand, plus la prédiction est stable & 50, 100 & 100\\
\texttt{max\_depth} & Profondeur maximale : limite le surapprentissage & 10, 15 & 15\\
\bottomrule
\end{tabularx}
\end{table}
Seule la configuration (100 ; 15) a été exécutée dans le notebook ; elle n'a pas été comparée aux autres par validation.

%==================== 8 PROTOCOLE ====================
\section{Protocole expérimental}\label{sec:proto}
\begin{table}[H]\centering\small
\caption{Découpage des données}
\begin{tabular}{@{}lll@{}}
\toprule
Ensemble & Période ou proportion & Observations\\\midrule
Apprentissage & 80\,\% premières lignes triées par année & 6\,630\\
Validation & non utilisée (pas de validation croisée) & ---\\
Test & 20\,\% dernières lignes (années les plus récentes) & 1\,658\\
\bottomrule
\end{tabular}
\end{table}
\begin{itemize}[leftmargin=1.2em]
\item \textbf{Référence naïve :} \texttt{DummyRegressor(strategy='mean')}, qui prédit la moyenne de la cible d'apprentissage.
\item \textbf{Validation :} holdout chronologique unique, sans réglage d'hyperparamètres.
\item \textbf{Métrique principale :} $R^2$, part de variance expliquée.
\item \textbf{Métriques secondaires :} MAE et RMSE, en USD.
\item \textbf{Graine aléatoire :} 42.
\end{itemize}

%==================== 9 RESULTATS ====================
\section{Résultats}\label{sec:res}
\begin{table}[H]\centering\small
\caption{Comparaison de la référence naïve et du modèle (jeu de test)}
\begin{tabular}{@{}lrrr@{}}
\toprule
Modèle & $R^2$ & MAE (USD) & RMSE (USD)\\\midrule
Référence : moyenne & $-0{,}0099$ & 208\,898\,312 & ---\\
Forêt aléatoire & $0{,}7321$ & 87\,935\,473 & 586\,000\,572\\
\bottomrule
\end{tabular}
\end{table}
\textbf{Lecture des résultats.} Le modèle explique 73,2\,\% de la variance, alors que la référence naïve a un $R^2$ légèrement négatif. La MAE passe de 208,9 à 87,9 millions USD, soit une baisse de 57,9\,\%.

\begin{figure}[H]\centering
\includegraphics[width=0.85\textwidth]{figures/fig4_reel_vs_predit.png}
\caption{Valeurs réelles contre valeurs prédites sur le jeu de test (millions USD)}\label{fig:pred}
\end{figure}
\textbf{Analyse des erreurs.} La plupart des points se regroupent près de la diagonale pour les petites et moyennes valeurs. Pour les très grosses transactions (au-delà de 8\,000 millions USD), le modèle sous-estime fortement : les prédictions plafonnent vers 7\,500 millions USD. C'est le comportement attendu d'une forêt, qui ne prédit pas au-delà des valeurs vues à l'entraînement. Ces quelques points pèsent lourd dans le RMSE (586 millions USD), très supérieur à la MAE.

%==================== 10 DISCUSSION ====================
\section{Discussion des résultats}\label{sec:disc}
\textbf{Réponse à l'objectif.} L'objectif est atteint : $R^2=0{,}7321>0{,}65$ et réduction de MAE de 57,9\,\%, supérieure au seuil de 30\,\%.

\textbf{Interprétation.} Le modèle fait mieux que la moyenne car partenaire et produit identifient la taille habituelle d'un flux. Aucune mesure d'importance des variables n'a été faite : on ne peut donc pas dire laquelle pèse le plus, ni conclure à une causalité.

\textbf{Surapprentissage et généralisation.} Le score d'apprentissage n'a pas été calculé et un seul découpage a été testé : l'écart apprentissage--test et la stabilité ne sont donc pas connus.

\textbf{Comparaison avec les travaux liés.} Les travaux de gravité \cite{anderson,yotov} s'appuient sur des variables macroéconomiques absentes ici ; les résultats ne sont pas directement comparables.

\textbf{Limites et biais.}
\begin{itemize}[leftmargin=1.2em]
\item Pas de PIB, taux de change ni coût de transport.
\item Le One-Hot sur \texttt{partnerDesc} crée beaucoup de colonnes ; \texttt{cmdCode} est traité comme un nombre alors que c'est une catégorie.
\item Les lignes « World » et les lignes par pays peuvent se recouvrir, ce qui fausse les sommes des figures~\ref{fig:evol} et~\ref{fig:part} et peut gonfler la cible.
\item Un seul holdout, sensible aux chocs d'une année (inflation après 2022).
\end{itemize}

\textbf{Pistes d'amélioration.}
\begin{enumerate}[leftmargin=1.5em]
\item Encodage cible ou regroupement des partenaires pour réduire la dimension, et traiter \texttt{cmdCode} comme catégorie.
\item Ajouter des données macroéconomiques externes (PIB, inflation du partenaire) et retirer les agrégats « World ».
\end{enumerate}

%==================== 11 CONCLUSION ====================
\section{Conclusion}\label{sec:concl}
Oui, la valeur des transactions d'exportation marocaines peut être prédite de façon utile à partir du partenaire, du produit, de l'année et du flux : la forêt aléatoire atteint un $R^2$ de 0,732 et une MAE de 87,9 millions USD, contre 208,9 pour la référence naïve. Ce résultat offre une base pour anticiper les rentrées de devises, mais il sous-estime les très gros flux. En perspective : nettoyer les agrégats, ajouter des variables macroéconomiques et valider avec une découpe temporelle répétée.

%==================== REFERENCES ====================
\begin{thebibliography}{9}
\bibitem{comtrade} United Nations Statistics Division, \emph{UN Comtrade Database}, \url{https://comtradeplus.un.org/}, extraction du 19/11/2026.
\bibitem{tinbergen} J. Tinbergen, \emph{Shaping the World Economy: Suggestions for an International Economic Policy}, Twentieth Century Fund, 1962.
\bibitem{anderson} J. E. Anderson et E. van Wincoop, « Gravity with Gravitas: A Solution to the Border Puzzle », \emph{American Economic Review}, vol.~93, n°1, p.~170--192, 2003.
\bibitem{yotov} Y. V. Yotov, R. Piermartini, J.-A. Monteiro et M. Larch, \emph{An Advanced Guide to Trade Policy Analysis: The Structural Gravity Model}, OMC, 2016.
\bibitem{breiman} L. Breiman, « Random Forests », \emph{Machine Learning}, vol.~45, n°1, p.~5--32, 2001.
\bibitem{sklearn} F. Pedregosa et al., « Scikit-learn : Machine Learning in Python », \emph{Journal of Machine Learning Research}, vol.~12, p.~2825--2830, 2011.
\bibitem{cours} Support du cours BD \& DS, notebook Tp\_ML.ipynb, ENCG-2026.
\end{thebibliography}

\appendix
\section{Reproduction}
Dépôt : \url{https://github.com/wardataisa/25007836-Taisa-CI.git}, étiquette \texttt{remise-finale}.
\begin{enumerate}[leftmargin=1.5em]
\item Cloner le dépôt : \texttt{git clone} suivi du lien ci-dessus.
\item Installer les dépendances : \texttt{pip install -r requirements.txt}.
\item Ouvrir le notebook \texttt{Modele\_Projet\_CI\_S7\_ENCG\_Settat.ipynb} puis « Redémarrer et tout exécuter ». Le fichier CSV Comtrade doit être placé dans le dossier indiqué dans la cellule de chargement.
\end{enumerate}

\section{Tableaux et figures complémentaires}
Aucun élément complémentaire.

\section{Déclaration des outils d'IA}
\begin{table}[H]\centering\small
\caption{Ressources et outils d'IA utilisés}
\begin{tabularx}{\textwidth}{@{}XXXX@{}}
\toprule
Outil ou ressource & Usage précis & Parties concernées & Vérification effectuée\\\midrule
Claude (Anthropic), 26/11/2026 & Mise en forme du rapport en LaTeX et rédaction du texte à partir du notebook & Rapport entier (sections 1 à 11, annexes) & Chiffres et figures recoupés avec les sorties du notebook ; sources à contrôler par l'étudiante\\
\bottomrule
\end{tabularx}
\end{table}
\end{document}
